% Input:       string declarations (waypoints, windings, detours, crossings)
%              from the language stage; frames and silhouettes from geometry.
% Output:      saved smooth paths, resolved crossing gaps, bead anchors, and
%              string events; development ink until the render stage owns it.
% Owned state: the string registry (ids, kinds, paths) and the crossing ledger.
% Invariants:  every declared crossing intersects; every intersection between
%              two strings is declared; paths exist before resolution.
% Next stage:  tenkz-render draws strings; the language stage supplies records.
%
% Curve model.  A string routes through waypoints as a Hobby spline: Hobby's
% algorithm (the METAFONT curve engine) chooses tangents and curvature
% globally, so a route declared by positions alone comes out graceful with
% no angle annotations.  This retires the out=/in= angle pair, the largest
% escape-hatch population in the shrink census (297 corpus occurrences).
% Closed strings are closed Hobby cycles: one smooth loop replaces the
% four-arc closure quartet.  Crossings are declared data: the under path is
% split at its geometric intersections with the over path and each cut end
% retreats by half the crossgap; the over path is drawn whole.  Storage and
% surgery are spath3; intersection counting is the pgf intersections engine;
% both are xelatex-clean and verified composable with Hobby splines here.
% The consumer loads the TikZ libraries hobby, spath3, and intersections.
%
% Units.  Waypoints and centres are caller-space TikZ numbers, written as
% braced pairs {x,y}; silhouette queries return pitch units and convert at
% the boundary through an explicit caller-space scale.  Every styling
% constant is a named metric ratio.

\makeatletter
\ExplSyntaxOn

% ---------- metric entries ---------------------------------------------------
% crossgap: total visual break in the under strand at a crossing, as a pitch
% ratio.  Calibration: quantikz opens 3pt at its reference line width; at the
% 11mm pitch that is 3pt/31.2pt = 0.096, rounded to one clean step.
\__tenkz_metric_declare:nnn {stringcrossgap} {0.10}
  { quantikz~opens~3pt~at~11mm~pitch;~the~community-calibrated~crossing~break }
% tube: meridian amplitude of a wound string, as a pitch ratio.  Large enough
% that one oscillation reads at a four-station ring, small enough that the
% winding never strays a full station gap from the cycle it decorates.
\__tenkz_metric_declare:nnn {stringtube} {0.35}
  { winding~oscillation~legible~at~four~stations~without~leaving~the~ring }

% ---------- diagnostics ------------------------------------------------------
\msg_new:nnnn {tenkz}{string-redeclared}
  { [TKZ-STRING-REDECLARED]~string~'#1'~already~exists }
  { Each~string~id~is~declared~once;~pick~a~fresh~name. }
\msg_new:nnnn {tenkz}{string-unknown}
  { [TKZ-STRING-UNKNOWN]~no~string~named~'#1' }
  { Declare~the~string~before~referencing~it. }
\msg_new:nnnn {tenkz}{string-nocross}
  { [TKZ-STRING-NOCROSS]~'#1'~is~declared~under~'#2'~but~they~never~meet }
  { The~declared~crossing~has~no~geometric~intersection;~fix~the~route~or~
    remove~the~declaration. }
\msg_new:nnnn {tenkz}{string-undeclared}
  { [TKZ-STRING-UNDECLARED]~strings~'#1'~and~'#2'~intersect~with~no~
    declared~order }
  { Every~crossing~carries~a~declared~over/under;~add~one~or~reroute. }
\msg_new:nnnn {tenkz}{string-library}
  { [TKZ-STRING-LIBRARY]~TikZ~library~'#1'~is~not~loaded }
  { The~string~engine~needs~hobby,~spath3,~and~intersections. }

% ---------- state ------------------------------------------------------------
\seq_new:N \l__tenkz_string_ids_seq          % declaration order
\prop_new:N \l__tenkz_string_kind_prop       % id -> open|closed|wind|around
\seq_new:N \l__tenkz_string_cross_seq        % {under}{over} pairs
\seq_new:N \l__tenkz_string_pts_seq
\seq_new:N \l__tenkz_string_scan_seq
\tl_new:N \l__tenkz_string_pts_tl
\tl_new:N \l__tenkz_string_first_tl
\tl_new:N \l__tenkz_string_last_tl
\tl_new:N \l__tenkz_string_tmp_tl
\tl_new:N \l__tenkz_string_id_tl
\fp_new:N \l__tenkz_string_x_fp
\fp_new:N \l__tenkz_string_y_fp
\fp_new:N \l__tenkz_string_dir_fp
\fp_new:N \l__tenkz_string_scale_fp
\int_new:N \l__tenkz_string_n_int
\bool_new:N \l__tenkz_string_hit_bool

\cs_new_protected:Npn \__tenkz_string_reset:
  {
    \seq_clear:N \l__tenkz_string_ids_seq
    \prop_clear:N \l__tenkz_string_kind_prop
    \seq_clear:N \l__tenkz_string_cross_seq
  }

\cs_new_protected:Npn \__tenkz_string_require_libs:
  {
    \cs_if_exist:cF {l__tikzspath_prefix_tl}
      { \msg_error:nnn {tenkz}{string-library} {spath3} }
    \cs_if_exist:cF {hobby@this@opts}
      { \msg_error:nnn {tenkz}{string-library} {hobby} }
    \cs_if_exist:NF \pgfintersectionofpaths
      { \msg_error:nnn {tenkz}{string-library} {intersections} }
  }

\cs_new_protected:Npn \__tenkz_string_register:nn #1#2
  {
    \prop_if_in:NnT \l__tenkz_string_kind_prop {#1}
      { \msg_error:nnn {tenkz}{string-redeclared} {#1} }
    \seq_put_right:Nn \l__tenkz_string_ids_seq {#1}
    \prop_put:Nnn \l__tenkz_string_kind_prop {#1} {#2}
  }

% load a token list of braced pairs {x,y}{x,y}... into the point seq
\cs_new_protected:Npn \__tenkz_string_load_pts:n #1
  {
    \seq_clear:N \l__tenkz_string_pts_seq
    \tl_map_inline:nn {#1}
      { \seq_put_right:Nn \l__tenkz_string_pts_seq {##1} }
  }

% ---------- path construction ------------------------------------------------
% An open string through braced pairs: first and last are endpoints, the
% interior are Hobby waypoints.  Two points degenerate to a straight bond,
% which the language stage forbids for kind=string; the engine tolerates it.
\cs_new_protected:Npn \__tenkz_string_declare_open:nn #1#2
  {
    \__tenkz_string_register:nn {#1} {open}
    \__tenkz_string_load_pts:n {#2}
    \int_set:Nn \l__tenkz_string_n_int
      { \seq_count:N \l__tenkz_string_pts_seq }
    \seq_pop_left:NN \l__tenkz_string_pts_seq \l__tenkz_string_first_tl
    \seq_pop_right:NN \l__tenkz_string_pts_seq \l__tenkz_string_last_tl
    \seq_set_map:NNn \l__tenkz_string_pts_seq \l__tenkz_string_pts_seq
      { ( ##1 ) }
    \tl_set:Ne \l__tenkz_string_pts_tl
      { \seq_use:Nn \l__tenkz_string_pts_seq { ~ } }
    \seq_if_empty:NTF \l__tenkz_string_pts_seq
      {
        \use:e
          {
            \exp_not:N \path [ spath/save = #1 ]
              ( \exp_not:o \l__tenkz_string_first_tl ) --
              ( \exp_not:o \l__tenkz_string_last_tl ) ;
          }
      }
      {
        \use:e
          {
            \exp_not:N \path [ spath/save = #1 ]
              ( \exp_not:o \l__tenkz_string_first_tl )
              to [ curve~through = { \exp_not:o \l__tenkz_string_pts_tl } ]
              ( \exp_not:o \l__tenkz_string_last_tl ) ;
          }
      }
    \tenkz@event{string|id=#1|kind=open|pts=\int_use:N \l__tenkz_string_n_int}
  }

% A closed string: one smooth Hobby cycle through every listed point.
\cs_new_protected:Npn \__tenkz_string_declare_closed:nn #1#2
  {
    \__tenkz_string_register:nn {#1} {closed}
    \__tenkz_string_load_pts:n {#2}
    \int_set:Nn \l__tenkz_string_n_int
      { \seq_count:N \l__tenkz_string_pts_seq }
    \seq_pop_left:NN \l__tenkz_string_pts_seq \l__tenkz_string_first_tl
    \seq_set_map:NNn \l__tenkz_string_pts_seq \l__tenkz_string_pts_seq
      { ( ##1 ) }
    \tl_set:Ne \l__tenkz_string_pts_tl
      { \seq_use:Nn \l__tenkz_string_pts_seq { ~ .. ~ } }
    \use:e
      {
        \exp_not:N \path [ spath/save = #1 , use~Hobby~shortcut ]
          ( [ closed ] \exp_not:o \l__tenkz_string_first_tl )
          .. ~ \exp_not:o \l__tenkz_string_pts_tl ;
      }
    \tenkz@event{string|id=#1|kind=closed|pts=\int_use:N \l__tenkz_string_n_int}
  }

% A wound string on a circle of centre (#2,#3) and radius #4 in caller space:
% homotopy class (p,q) = (#5,#6), meridian amplitude #7 in caller space
% (pass 0 for a pure longitude).  The radius is modulated as
% r + a cos(360 q t) while the angle sweeps 360 p t -- the planar projection
% of a torus winding; samples are dense enough that the closed Hobby cycle
% is indistinguishable from the parametric curve.
\cs_new_protected:Npn \__tenkz_string_declare_wind:nnnnnnn #1#2#3#4#5#6#7
  {
    \__tenkz_string_register:nn {#1} {wind}
    \int_set:Nn \l__tenkz_string_n_int
      { \fp_eval:n { max( 12 , 8 * max( #5 , #6 , 1 ) ) } }
    \tl_clear:N \l__tenkz_string_pts_tl
    \int_step_inline:nnn { 1 } { \l__tenkz_string_n_int - 1 }
      {
        \tl_put_right:Ne \l__tenkz_string_pts_tl
          { .. ~ ( \__tenkz_string_wind_pt:nnnnnnn
              {##1} {#2}{#3}{#4}{#5}{#6}{#7} ) ~ }
      }
    \use:e
      {
        \exp_not:N \path [ spath/save = #1 , use~Hobby~shortcut ]
          ( [ closed ]
            \__tenkz_string_wind_pt:nnnnnnn {0} {#2}{#3}{#4}{#5}{#6}{#7} )
          \exp_not:o \l__tenkz_string_pts_tl ;
      }
    \tenkz@event{string|id=#1|kind=wind|class=#5,#6|
      pts=\int_use:N \l__tenkz_string_n_int}
  }
% sample i of n: angle 360 p i/n, radius r + a cos(360 q i/n)
\cs_new:Npn \__tenkz_string_wind_pt:nnnnnnn #1#2#3#4#5#6#7
  {
    \fp_eval:n
      { round( #2 + (#4 + #7 * cosd( 360 * #6 * #1 / \l__tenkz_string_n_int ))
               * cosd( 360 * #5 * #1 / \l__tenkz_string_n_int ) , 5 ) }
    ,
    \fp_eval:n
      { round( #3 + (#4 + #7 * cosd( 360 * #6 * #1 / \l__tenkz_string_n_int ))
               * sind( 360 * #5 * #1 / \l__tenkz_string_n_int ) , 5 ) }
  }

% A detour past one measured obstacle: string #1 from point #2 to point #3
% (braced pairs), apex pushed along direction angle #4 to the obstacle's
% support plus daylight.  #5 is the obstacle as {family}{hx}{hy}{corner}
% {theta}{cx}{cy} in pitch units; #6 is the caller-space length of one pitch
% unit.  Geometry's support functions are consumed read-only.
\cs_new_protected:Npn \__tenkz_string_declare_around:nnnnnn #1#2#3#4#5#6
  {
    \__tenkz_string_register:nn {#1} {around}
    \tl_set:Nn \l__tenkz_string_id_tl {#1}
    \fp_set:Nn \l__tenkz_string_dir_fp {#4}
    \fp_set:Nn \l__tenkz_string_scale_fp {#6}
    \__tenkz_string_apex:nnnnnnn #5
    \use:e
      {
        \exp_not:N \path [ spath/save = #1 ]
          ( \exp_not:n {#2} )
          to [ curve~through =
            { ( \fp_use:N \l__tenkz_string_x_fp ,
                \fp_use:N \l__tenkz_string_y_fp ) } ]
          ( \exp_not:n {#3} ) ;
      }
    \tenkz@event{string|id=#1|kind=around|
      apexx=\fp_eval:n { round( \l__tenkz_string_x_fp , 5 ) }|
      apexy=\fp_eval:n { round( \l__tenkz_string_y_fp , 5 ) }}
  }
% apex = centre + (support(d) + daylight) * unit(d), pitch units scaled out
\cs_new_protected:Npn \__tenkz_string_apex:nnnnnnn #1#2#3#4#5#6#7
  {
    \fp_set:Nn \l__tenkz_string_x_fp
      { ( #6 + ( \__tenkz_geom_support:nnnnnn
                   {#1}{#2}{#3}{#4}{#5}{ \l__tenkz_string_dir_fp }
                 + \__tenkz_string_ratio:n {daylight} )
               * cosd( \l__tenkz_string_dir_fp ) )
        * \l__tenkz_string_scale_fp }
    \fp_set:Nn \l__tenkz_string_y_fp
      { ( #7 + ( \__tenkz_geom_support:nnnnnn
                   {#1}{#2}{#3}{#4}{#5}{ \l__tenkz_string_dir_fp }
                 + \__tenkz_string_ratio:n {daylight} )
               * sind( \l__tenkz_string_dir_fp ) )
        * \l__tenkz_string_scale_fp }
  }
% a metric ratio as a plain factor, for pitch-unit arithmetic
\cs_new:Npn \__tenkz_string_ratio:n #1
  { \prop_item:Nn \g__tenkz_metric_value_prop {#1} }

% ---------- crossings ----------------------------------------------------------
\cs_new_protected:Npn \__tenkz_string_under:nn #1#2
  {
    \prop_if_in:NnF \l__tenkz_string_kind_prop {#1}
      { \msg_error:nnn {tenkz}{string-unknown} {#1} }
    \prop_if_in:NnF \l__tenkz_string_kind_prop {#2}
      { \msg_error:nnn {tenkz}{string-unknown} {#2} }
    \seq_put_right:Nn \l__tenkz_string_cross_seq { {#1} {#2} }
  }

% intersection count of two saved paths through the pgf engine
\cs_new_protected:Npn \__tenkz_string_count:nnN #1#2#3
  {
    \pgfintersectionofpaths
      { \exp_args:Nc \pgfsetpath { \l__tikzspath_prefix_tl #1 } }
      { \exp_args:Nc \pgfsetpath { \l__tikzspath_prefix_tl #2 } }
    \int_set:Nn #3 { \pgfintersectionsolutions }
  }

% Resolve every declared crossing, then police the honesty rule: an
% intersecting undeclared pair is a hard error.  Runs inside the picture,
% after every declaration and before any string is drawn.
\cs_new_protected:Npn \__tenkz_string_resolve:
  {
    \seq_map_inline:Nn \l__tenkz_string_cross_seq
      { \__tenkz_string_resolve_one:nn ##1 }
    \__tenkz_string_police:
  }
\cs_new_protected:Npn \__tenkz_string_resolve_one:nn #1#2
  {
    \__tenkz_string_count:nnN {#1} {#2} \l__tenkz_string_n_int
    \int_compare:nNnTF { \l__tenkz_string_n_int } = { 0 }
      { \msg_error:nnnn {tenkz}{string-nocross} {#1} {#2} }
      {
        \tikzset
          {
            spath/split~at~intersections~with = {#1} {#2} ,
            spath/insert~gaps~after~components =
              {#1} { \dim_eval:n { \__tenkz_dim:n {stringcrossgap} } } ,
          }
        \tenkz@event{stringcross|under=#1|over=#2|
          hits=\int_use:N \l__tenkz_string_n_int}
      }
  }
\cs_new_protected:Npn \__tenkz_string_police:
  {
    \seq_set_eq:NN \l__tenkz_string_scan_seq \l__tenkz_string_ids_seq
    \seq_map_inline:Nn \l__tenkz_string_ids_seq
      {
        \seq_pop_left:NN \l__tenkz_string_scan_seq \l__tenkz_string_tmp_tl
        \seq_map_inline:Nn \l__tenkz_string_scan_seq
          { \__tenkz_string_police_pair:nn {##1} {####1} }
      }
  }
\cs_new_protected:Npn \__tenkz_string_police_pair:nn #1#2
  {
    \bool_set_false:N \l__tenkz_string_hit_bool
    \seq_map_inline:Nn \l__tenkz_string_cross_seq
      {
        \tl_if_eq:nnT { {#1} {#2} } {##1}
          { \bool_set_true:N \l__tenkz_string_hit_bool }
        \tl_if_eq:nnT { {#2} {#1} } {##1}
          { \bool_set_true:N \l__tenkz_string_hit_bool }
      }
    \bool_if:NF \l__tenkz_string_hit_bool
      {
        \__tenkz_string_count:nnN {#1} {#2} \l__tenkz_string_n_int
        \int_compare:nNnT { \l__tenkz_string_n_int } > { 0 }
          { \msg_error:nnnn {tenkz}{string-undeclared} {#1} {#2} }
      }
  }

% ---------- beads ---------------------------------------------------------------
% The point a fraction #2 along string #1, exposed as TikZ coordinate #3 and
% reported as an event.  The parameter is the spath segment parameter, which
% Hobby's near-uniform waypoint spacing keeps close to arclength; the event
% carries the resolved position so a fixture pins the truth, not the promise.
\cs_new_protected:Npn \__tenkz_string_bead:nnn #1#2#3
  {
    \prop_if_in:NnF \l__tenkz_string_kind_prop {#1}
      { \msg_error:nnn {tenkz}{string-unknown} {#1} }
    \tenkz@string@beadpath {#1} {#2} {#3}
    \pgfpointanchor {#3} {center}
    \tenkz@event{stringbead|id=#1|t=#2|
      x=\fp_eval:n { round( \dim_to_fp:n { \pgf@x } , 4 ) }pt|
      y=\fp_eval:n { round( \dim_to_fp:n { \pgf@y } , 4 ) }pt}
  }

% ---------- development ink -----------------------------------------------------
% The render stage owns string ink at integration; until then fixtures draw
% through this single, clearly marked primitive.
\cs_new_protected:Npn \__tenkz_string_dev_draw:nn #1#2
  {
    \prop_if_in:NnF \l__tenkz_string_kind_prop {#1}
      { \msg_error:nnn {tenkz}{string-unknown} {#1} }
    \draw [ spath/use = #1 , #2 ] ;
  }

\ExplSyntaxOff

% The spath coordinate system is parsed with document catcodes; this helper
% stays outside the expl3 region so 'spath cs:' reaches TikZ verbatim.
\def\tenkz@string@beadpath#1#2#3{%
  \path (spath cs:#1 #2) coordinate (#3);%
}

\makeatother
\endinput
